\documentclass[11pt]{article}

\usepackage[margin=1.02in]{geometry}
\usepackage[T1]{fontenc}
\usepackage{lmodern}
\usepackage{amsmath,amssymb,amsthm,mathtools}
\usepackage{array,booktabs,longtable}
\usepackage{xcolor}
\usepackage{microtype}
\usepackage{hyperref}
\hypersetup{
  colorlinks=true,
  linkcolor=blue!55!black,
  citecolor=blue!55!black,
  urlcolor=blue!55!black,
  pdftitle={Exclusion of the fixed-quadratic primitive-pencil quartic row}
}

\newtheorem{theorem}{Theorem}
\newtheorem{proposition}[theorem]{Proposition}
\newtheorem{lemma}[theorem]{Lemma}
\newtheorem{corollary}[theorem]{Corollary}
\theoremstyle{remark}
\newtheorem{remark}[theorem]{Remark}

\newcommand{\A}{\mathbb A}
\newcommand{\C}{\mathbb C}
\newcommand{\PP}{\mathbb P}
\newcommand{\Jac}{\operatorname{Jac}}
\newcommand{\rank}{\operatorname{rank}}
\newcommand{\ord}{\operatorname{ord}}

\title{\bfseries Exclusion of the fixed-quadratic primitive-pencil
quartic row\\[3pt]
\large A frozen-case structural theorem for Keller maps in dimension three}
\author{Alec Kriebel\\with substantial AI assistance}
\date{First public repository release (UTC):
  \texttt{\detokenize{2026-07-26T00:45:00Z}}}

\begin{document}
\maketitle

\begin{center}
\fcolorbox{black!45}{yellow!8}{%
\begin{minipage}{0.91\textwidth}
\small
\textbf{Review, priority, and verification notice.}
This manuscript is not peer reviewed.  Exact computer algebra checks the
identities and finite case ledgers encoded in the accompanying repository;
it is evidence, not peer review.  Substantial AI assistance was used in the
mathematical exploration, proof organization, software, hostile internal
audits, and exposition.  No claim of worldwide priority is made.  The result
excludes one of fourteen frozen quartic leading-form rows; it does not prove
the full degree-four case.
\end{minipage}}
\end{center}

\begin{abstract}
Let $F:\A^3_\C\to\A^3_\C$ be a polynomial map of total degree four with
nonzero constant Jacobian.  We prove that $F$ is an automorphism when its
quartic homogeneous part has a fixed quadratic component and a primitive
quadratic line pencil.  In the program's canonical notation, the excluded
row is
\[
(\rank JH_4,e,a,b,\delta,\nu)=(2,2,2,1,1,1).
\]
The proof first gives a coefficient-independent normalization
$H_4=(hp,hq,0)$ and uses valuations of homogeneous first integrals to reduce
to a unique-double-line frontier.  Keeping separate the fixed quadratic
$h$, the double member $s$, and the cubic companion produces a genuine
$3+\PP^1+3$ orbit space, rather than a finite companion list.  Exact
weighted-Jacobian identities exclude the whole projective family, its six
endpoints, and two discrete open orbits.  A clean-room dependency-free
reconstruction verifies the frozen coefficient routing and the two discrete
calculations.  The universal certified total-degree floor remains four.
\end{abstract}

\section{Invariant statement}

Translate the source and target so that
\begin{equation}
F=LX+H_2+H_3+H_4,\qquad L\in\operatorname{GL}_3(\C),
\label{eq:F}
\end{equation}
where $H_i$ is homogeneous of degree $i$.  Suppose that
$\rank JH_4=2$.  Extract
\[
h=\gcd(H_{4,1},H_{4,2},H_{4,3}),\qquad e=\deg h,
\qquad G=H_4/h.
\]
Let $E_G$ be the relative algebraic closure of
$\C(G_i/G_j)$ in $\C(\PP^2)$.  It is rational, so
$E_G=\C(p/q)$ for a coprime homogeneous pair of least common degree $a$.
If the induced binary target triple has degree $b$, the reduced image has
degree $\delta$, and the map to its normalization has degree $\nu$, then
\[
e+ab=4,\qquad b=\delta\nu.
\]

\begin{theorem}[Fixed-quadratic primitive-pencil exclusion]
\label{thm:main}
Let $F$ be as in \eqref{eq:F}, over $\C$, of exact total degree four.
Suppose
\[
(\rank JH_4,e,a,b,\delta,\nu)=(2,2,2,1,1,1).
\]
If $F$ has nonzero constant Jacobian determinant, then $F$ is a polynomial
automorphism.
\end{theorem}

\begin{corollary}
No degree-four Keller counterexample belongs to frozen row
\texttt{Q2-E2-A2-B1-D1-N1}.
\end{corollary}

The fourteen-row frozen denominator currently has
\[
4/14\ \text{certified excluded},\qquad
4/14\ \text{provisional},\qquad
6/14\ \text{open}.
\]
This bookkeeping is not a statement about the other rows.

\section{Normalization and weighted identities}

\begin{proposition}\label{prop:normal}
After an invertible target change,
\begin{equation}
H_4=(hp,hq,0),                                      \label{eq:normal}
\end{equation}
where $h,p,q$ are quadratic, $\gcd(p,q)=1$, and
$\C(p/q)$ is relatively algebraically closed in $\C(\PP^2)$.
Every frozen coefficient chart is covered without dividing by its first
nonzero quartic coefficient.
\end{proposition}

\begin{proof}
The canonical factorization is $H_4=hA(p,q)$, where
$A(u,v)=\mathbf a u+\mathbf b v$ and $\mathbf a,\mathbf b\in\C^3$ are
independent.  Complete them to a target basis.  This gives
\eqref{eq:normal} and transports all lower jets bijectively.

For a division-free coverage statement, write the coefficients of the
primitive quadratic triple as a $3\times6$ matrix $B$.  Its rank is two.
Order its $45=3\binom62$ minors and choose the first nonzero one.  The
corresponding two target rows span the third, and the cleared row relation
has determinant equal, up to sign, to the selected minor.  Thus it gives
\eqref{eq:normal}.  Among the frozen first-nonzero coefficient guards,
\texttt{C00}--\texttt{C29} reach this minor router.  The guards
\texttt{C30}--\texttt{C44} make the first two target components zero, so
$\rank JH_4\le1$; they are empty in the present row.
\end{proof}

Put
\[
\mathcal J(w)=L+wJH_2+w^2JH_3+w^3JH_4,\qquad
E_j=[w^j]\det\mathcal J(w).
\]
The Keller condition gives
\begin{equation}
E_9=\cdots=E_1=0,\qquad\det L\ne0.                  \label{eq:weighted}
\end{equation}
Writing $P=hp$, $Q=hq$, and $R=(H_3)_3$, the zero third row of $JH_4$
gives
\begin{equation}
E_8=\Jac(P,Q,R)=0.                                  \label{eq:E8}
\end{equation}

\section{The top-identity frontier}

\begin{lemma}[First-integral descent]\label{lem:descent}
If $0\ne T$ is homogeneous of degree $d$ and
$\Jac(P,Q,T)=0$, then
\[
\frac{T^4}{P^d}=\mathcal R(q/p)
\]
for some $\mathcal R\in\C(t)$.
\end{lemma}

\begin{proof}
The Jacobian identity makes $T$ algebraic over $\C(P,Q)$.  The displayed
quotient has degree zero, hence lies in $\C(\PP^2)$.  Separating the scaling
coordinate shows that it is algebraic over $\C(q/p)$.  Relative algebraic
closure of the minimal pencil field then places it in $\C(q/p)$.
\end{proof}

A prime factor $f$ of $h$ is called horizontal if it is not contained in a
pencil fibre.  Taking $v_f$ in Lemma~\ref{lem:descent}, with $d=3$, gives
\[
4v_f(R)=3v_f(h).
\]
Since $v_f(h)$ is one or two, this is impossible for nonzero $R$.  Thus a
horizontal component forces $R=0$.

For the all-vertical case, comparing two reduced components of one pencil
fibre removes the unknown order of $\mathcal R$.  This yields the following
exact reduction.

\begin{proposition}[All-vertical reduction]\label{prop:vertical}
If every component of $h$ is vertical and \eqref{eq:E8} has a nonzero
cubic solution, then, after source and pencil changes,
\begin{equation}
H_4=(h^2,hs,0),\qquad s=\ell^2,\qquad
R=\ell r,\quad [r]\in\PP\langle h,s\rangle,          \label{eq:frontier}
\end{equation}
where $s$ is the pencil's unique double-line member.
The underlying pencil is one of
\[
\langle x^2,yz\rangle,\qquad
\langle x^2,y^2+xz\rangle.
\]
\end{proposition}

\begin{proof}
If $h=\ell^2$ and a vertical member is $\ell m$ with
$m\not\sim\ell$, comparison of the two components of that fibre gives
$4N=6$.  If $h=\ell_1\ell_2$ and the vertical members are distinct, the
same argument gives $4N=3$ unless both members are squares; the latter
would make $p/q$ a quadratic function of a linear pencil and contradict
minimality.

The remaining shape has one pencil member equal to $h$.  Then
\[
\Jac(h^2,hq,R)=2h^2\Jac(h,q,R).
\]
First-integral descent gives $R^2/h^3=S(q/h)$.  If every conic fibre were
reduced, all finite orders of $S$ would be even and the order at the
$h$-fibre odd, contradicting degree zero of a principal divisor on
$\PP^1$.  Hence the pencil contains a double line $s=\ell^2$.
Two distinct double members would again violate minimality, so $s$ is
unique.  Direct calculation gives the complete cubic kernel
$\ell\langle h,s\rangle$.

Put $s=x^2$.  The restriction of another member to $x=0$ has rank two or
one.  Completing cross terms gives respectively
$\langle x^2,yz\rangle$ and
$\langle x^2,y^2+xz\rangle$.  This proves the claim.
\end{proof}

\begin{lemma}[Quadratic-component exit]\label{lem:exit}
If a degree-at-most-four Keller map of $\A^3_\C$ has a nonzero
target-linear component of degree at most two, it is an automorphism.
\end{lemma}

\begin{proof}
A quadratic polynomial with nowhere-vanishing gradient can be straightened
to a coordinate by a polynomial automorphism of degree at most two.
The transformed map is $(F_1,F_2,z)$.  On every $z$-fibre the first two
components form a plane etale map of degree at most eight.  Moh's
unconditional bounded plane theorem, in the degree-at-most-twelve form
quoted by Vistoli, makes every fibre map an automorphism
\cite{Moh,Vistoli}.  Hence the threefold map is injective and therefore an
automorphism.
\end{proof}

Consequently every branch with $R=0$ exits the counterexample locus by
Lemma~\ref{lem:exit}.  It is important that this is an automorphism
conclusion, not a contradiction to the existence of a constant Jacobian.

\section{The marked companion quotient}

The frontier \eqref{eq:frontier} has three pieces of intrinsic data:
the pencil, its double member $[s]$, the fixed gcd $[h]$, and the companion
$[r]$.  The loci $h=s$ and $h\ne s$ must be treated separately.

When $h=s$, the two canonical pencils above have exactly two nonzero
companion orbits:
\[
R=x^3,\qquad R=xq.
\]
Exact weighted-determinant eliminations exclude both orbits for both
pencils.  These four calculations are the marked-equal ancillary packages.

For $h\ne s$, the marked pair and companion have exactly the following
thirteen stable strata.

\begin{longtable}{lll}
\toprule
marked pair $(s,h)$ & companion suffix & representative or guard\\
\midrule
\endhead
$(x^2,yz)$ & C0 & $R=0$\\
            & CH & $r=h$\\
            & CS & $r=s$\\
            & CO & $r\notin\{h,s\}$\\[2pt]
$(x^2,x^2+yz)$ & C0 & $R=0$\\
            & CH & $r=h$\\
            & CT & $r=yz$\\
            & CS & $r=s$\\
            & CTAU & $r=h+ks,\ k\ne0,-1$\\[2pt]
$(x^2,y^2+xz)$ & C0 & $R=0$\\
            & CH & $r=h$\\
            & CS & $r=s$\\
            & CO & $r\notin\{h,s\}$\\
\bottomrule
\end{longtable}

The middle stabilizer fixes its pencil line pointwise, so CTAU is a genuine
one-parameter orbit family.  The projective values $0,-1,\infty$ are,
respectively, CH, CT, and CS; none may be removed by a generic calculation.
Two independent stabilizer reconstructions certify the count
\[
4+5+4=13,\qquad\text{nonzero orbit-space shorthand }3+\PP^1+3.
\]

\section{Uniform exclusion of the projective family}

On the finite chart put
\[
h=x^2+yz,\quad s=x^2,\quad r=h+ks,\quad k\ne0.
\]
The complete $E_7$ normal form, modulo two target shears and three source
translations, is
\[
H_3=(Ax^3,Bx^3,x(h+ks)),\qquad (H_2)_3=Tx^2.
\]
Two $E_7$ pivots cover every finite $k$ because
\[
\frac12(9k^2+6k-1)-\frac32(k+1)(3k-1)=1.
\]

Write the first two quadratic components in monomial order
$(x^2,xy,xz,y^2,yz,z^2)$ with coefficients $a_i,b_i$, and put
$L=(\ell_0,\ldots,\ell_8)$ by rows.  The $E_6$ chains give
\[
\begin{aligned}
b_1&=2\ell_7,&a_1&=-12k\ell_7,&-36k^2\ell_7&=0,\\
b_2&=2\ell_8,&a_2&=-12k\ell_8,& 36k^2\ell_8&=0,
\end{aligned}
\]
together with $a_3=a_5=b_3=b_5=0$.  Since $k\ne0$,
\[
a_1=a_2=a_3=a_5=b_1=b_2=b_3=b_5=\ell_7=\ell_8=0.
\]
Six literal $E_5$ coefficients then include
\[
-2\ell_4,\quad2\ell_5,\quad
-\ell_1-(6k+4)\ell_4,\quad
\ell_2+(6k+4)\ell_5,
\]
so $\ell_1=\ell_2=\ell_4=\ell_5=0$.  The last two columns of $L$
vanish, contrary to \eqref{eq:weighted}.  This uniformly excludes CTAU
and the finite CT boundary $k=-1$.

\section{Six endpoints and two discrete orbits}

The six CH/CS endpoint normal forms have four $E_6$ compatibility ideals:
\[
\begin{array}{c|l}
\text{type}&I_6\\ \hline
\mathrm{RT/H}&(AC,AD,AE,AF,CE,DF,E^2,F^2)\\
\mathrm{RT/S}&(C^2,D^2)\\
\mathrm{RO/H}&(AC,AD,AE,AF,CF+DE,EF,2DF-E^2,F^2)\\
\mathrm{RO/S}&(CD,D^2).
\end{array}
\]
Global $E_5$ syzygies and fresh solves on every pivot divisor force a
singular $L$ in five endpoint slices.  The sixth, RO/H, has a genuine
invertible survivor through $E_5$ on
\[
A=C=D=E=F=0,\qquad T\ne0.
\]
After the complete $E_6/E_5$ solve, two $E_4$ coefficients are
\[
[xyz^2]E_4=-8\ell_8^2,\qquad
[x^2yz]E_4=-4\bigl(2(b_0-\ell_6)\ell_8-\ell_7^2\bigr).
\]
They give $\ell_8=\ell_7=0$, while $\det L$ is divisible by $\ell_7$.
Thus all six endpoints close.

It remains to record the two outer CO calculations.  For
\[
(h,r)=(yz,x^2+yz)
\]
the raw $E_7$ matrix has rank $18$, with five legal gauge directions and
complete normal complement
\[
(x^3,0,0),\quad(0,x^3,0),\quad(0,0,x^2).
\]
The $E_6$ equations force
\[
a_1=a_2=a_3=a_5=b_1=b_2=b_3=b_5=\ell_7=\ell_8=0,
\]
and $E_5$ forces
$\ell_1=\ell_2=\ell_4=\ell_5=0$.

For
\[
(h,r)=(y^2+xz,x^2+y^2+xz)
\]
the complete normal complement is
\[
(x^3,0,0),\quad(0,x^3,0),\quad
(2z(y^2+xz),x^2z,xz).
\]
If $C$ is the last normal parameter, the $E_6$ equations give
\[
a_1=a_4=b_1=b_4=b_5=\ell_7=\ell_8=0,\quad
a_2=a_3,\quad a_5=C^2,\quad b_2=b_3.
\]
Then $E_5$ gives
\[
\ell_1=\ell_4=0,\qquad
\ell_2=Ca_3,\qquad\ell_5=Cb_3.
\]
The second column of $L$ is zero.  This calculation includes $C=0$ and
uses no localization.

\section{Assembly of the proof}

\begin{proof}[Proof of Theorem~\ref{thm:main}]
Proposition~\ref{prop:normal} covers every frozen coefficient guard.
If $h$ has a horizontal component, the valuation congruence forces $R=0$,
and Lemma~\ref{lem:exit} gives an automorphism.  If every component is
vertical, Proposition~\ref{prop:vertical} either gives the same exit or
reaches \eqref{eq:frontier}.

On the frontier, $R=0$ again gives an automorphism.  For $R\ne0$, the
marked-equal packages or exactly one of the thirteen marked-distinct
strata applies.  The four marked-equal nonzero orbits, the uniform
projective family, all six endpoints, and both discrete CO orbits force
$\det L=0$.  This contradicts the constant term of the Keller determinant.
No Keller counterexample remains in the row.
\end{proof}

\section{Verification and audit}

The aggregate command
\begin{verbatim}
./verify_full_row_strict.sh
\end{verbatim}
runs exact SymPy, independent PARI/GP, and dependency-free checks for the
top reduction, all lower packages, and the post-freeze bridge.  Its final
marker is
\begin{verbatim}
Q2_E2_A2_B1_D1_N1_FULL_ROW_STRICT_PASS_4D95A1
\end{verbatim}

The clean-room bridge auditor did not read the candidate bridge, the root
exploration scripts, or the primary CO and endpoint verifiers.  Its
standard-library sparse-polynomial engine independently reconstructs:
\[
30\text{ rank-two-possible frozen pivots},\quad
15\text{ empty pivots},\quad45\text{ intrinsic minor charts},
\]
the $4+5+4$ marked ledger, both raw CO quotients, their lower exits, and the
uniform finite chart.  Its marker is
\begin{verbatim}
AUDIT_BRIDGE_Q2_E2_STRICT_PASS_D9347B
\end{verbatim}

The audit found and corrected two draft errors: C0 is an automorphism exit
rather than a blanket determinant contradiction, and
\texttt{C30}--\texttt{C44} are empty rather than entering the minor router.
Exact checks verify the encoded algebra; they are not peer review.

\begin{thebibliography}{9}

\bibitem{Moh}
T.-T. Moh,
\emph{On the Jacobian conjecture and the configurations of roots},
J. Reine Angew. Math. \textbf{340} (1983), 140--212.

\bibitem{Vistoli}
Angelo Vistoli,
\emph{The Jacobian conjecture in dimension 3 and degree 3},
J. Pure Appl. Algebra \textbf{142} (1999), 79--89.

\end{thebibliography}

\end{document}
